\chapter{Unicompartmental Knee Replacement}

\section*{Introduction \& Clinical Indications}
Unicompartmental knee replacement (UKR), also known as partial knee replacement, is a minimally invasive surgical procedure aimed at resurfacing only the damaged compartment of the knee joint, most frequently the medial compartment. This technique is particularly beneficial for patients suffering from localized osteoarthritis, where degeneration is confined to a single compartment, allowing for the preservation of healthy cartilage, bone, and ligaments in the unaffected areas. Unlike total knee replacement, which involves the complete resurfacing of the knee joint, UKR maintains the integrity of the anterior and posterior cruciate ligaments, thereby enhancing postoperative function and stability. The primary goals of this surgery are to alleviate pain, restore mobility, and improve the overall quality of life for patients who meet specific criteria, making it a preferred option for those with isolated knee joint degeneration.

\begin{itemize}
  \item Partial Knee Replacement
  \item Unicondylar Arthroplasty
  \item Oxford Knee (referring to a specific implant design)
  \item Medial Unicompartmental Knee Arthroplasty (if medial compartment is replaced)
  \item Lateral Unicompartmental Knee Arthroplasty (if lateral compartment is replaced)
\end{itemize}

The primary surgical principle of unicompartmental knee replacement is to selectively excise the diseased portion of the knee joint while preserving healthy surrounding tissues. In cases of localized osteoarthritis, the degeneration of cartilage leads to pain and functional limitations due to bone-on-bone contact. The UKR procedure involves the removal of the damaged articular cartilage and a minimal amount of underlying bone from the affected compartment, which is then replaced with prosthetic components designed to restore joint function.

The rationale behind this approach is to maintain the natural biomechanics of the knee. By preserving the cruciate ligaments and the healthy compartment, the procedure aims to retain proprioception and kinematic pathways, resulting in a more natural-feeling knee postoperatively. Technical goals include achieving proper limb alignment, restoring joint line height, ensuring adequate soft tissue balance, and securing stable fixation of the prosthetic components to the bone. This targeted intervention not only alleviates pain but also enhances the patient's functional capabilities, allowing for a quicker return to daily activities.

The evolution of unicompartmental knee replacement began in the 1970s when surgeons sought less invasive alternatives to total knee arthroplasty. In 1972, Dr. Marmor introduced the first modern unicondylar implants, demonstrating the feasibility of resurfacing only one compartment of the knee. However, early designs faced challenges, including high rates of loosening and revision due to inadequate materials and fixation techniques.

Significant advancements occurred in the 1980s and 1990s, particularly with the development of the Oxford Knee system by Dr. Goodfellow and Dr. O'Connor, which featured a mobile bearing polyethylene insert. This innovation aimed to reduce wear and improve kinematics, leading to more durable outcomes. The refinement of surgical techniques, including smaller incisions and precise instrumentation, further enhanced the procedure's efficacy. In the 21st century, the introduction of computer-assisted navigation and robotic surgery has improved implant positioning accuracy, contributing to better outcomes and expanding the indications for UKR. These advancements have solidified UKR's role as a viable treatment option for patients with single-compartment knee osteoarthritis.

Unicompartmental knee replacement is indicated for several specific clinical scenarios, including:

\begin{itemize}
  \item Osteoarthritis confined to a single compartment of the knee, confirmed by imaging studies such as X-rays and MRI.
  \item Intact anterior and posterior cruciate ligaments, essential for knee stability post-UKR.
  \item Minimal or no arthritis in the patellofemoral joint or the contralateral tibiofemoral compartment.
  \item Correctable varus (bow-legged) or valgus (knock-kneed) deformity, typically less than 10-15 degrees.
  \item Patients who are generally active, of reasonable weight, and motivated for rehabilitation.
  \item Persistent knee pain and functional limitation despite non-operative treatments such as physical therapy, anti-inflammatory medications, and injections.
  \item Age is less of a strict contraindication than activity level and bone quality, though younger, active patients are often ideal candidates.
  \item Failure of previous arthroscopic procedures, such as debridement or meniscectomy, to provide lasting pain relief.
\end{itemize}

These indications highlight the importance of careful patient selection to ensure optimal outcomes following UKR.

Unicompartmental knee replacement is primarily considered a first-line surgical treatment for patients with end-stage osteoarthritis confined to a single compartment. It is not typically viewed as a secondary or salvage procedure, as its success relies on the preservation of healthy surrounding structures. For patients who meet the strict criteria, UKR offers distinct advantages over total knee replacement, including a smaller incision, less bone resection, faster recovery, and potentially better range of motion and proprioception. However, if the arthritis progresses to other compartments or if the patient's ligaments are compromised, UKR becomes inappropriate, and total knee replacement is then considered the primary treatment. In cases of UKR failure due to implant loosening or progression of arthritis, conversion to total knee replacement is often the secondary surgical solution.

\section*{Relevant Surgical Anatomy}
The knee joint is a complex structure composed of the femur, tibia, and patella, with the medial and lateral compartments separated by the cruciate ligaments. The medial compartment consists of the medial femoral condyle and the medial tibial plateau, while the lateral compartment includes the lateral femoral condyle and the lateral tibial plateau. The arterial supply to the knee is primarily provided by the popliteal artery, which branches into the anterior and posterior tibial arteries, while venous drainage occurs through the popliteal vein and its tributaries. 

Key anatomical landmarks include:

\begin{itemize}
  \item **McBurney's Point:** A reference point for the medial compartment.
  \item **Calot's Triangle:** Important for understanding the vascular supply and potential surgical approaches.
  \item **Cruciate Ligaments:** Anterior and posterior cruciate ligaments are vital for knee stability and must be preserved during UKR.
  \item **Menisci:** The medial and lateral menisci provide cushioning and stability to the knee joint.
\end{itemize}

Understanding these anatomical relationships is crucial for the successful execution of unicompartmental knee replacement, as it allows for precise surgical intervention while minimizing damage to surrounding structures.

\section*{Preoperative Workup \& Preparation}
Thorough preoperative preparation is essential for a successful unicompartmental knee replacement. This typically begins with a comprehensive medical evaluation to assess the patient's overall health and fitness for surgery.

\begin{itemize}
  \item **Medical Clearance:** Patients undergo a full physical examination, blood tests (complete blood count, metabolic panel, coagulation studies), and often an electrocardiogram (ECG) and chest X-ray. Cardiac and pulmonary clearance may be required from specialists, especially for older patients or those with pre-existing conditions.
  
  \item **Imaging Studies:** Standard weight-bearing X-rays of the knee (AP, lateral, skyline, and long-leg alignment views) are crucial to confirm single-compartment disease and assess overall limb alignment. A magnetic resonance imaging (MRI) scan is often performed to further evaluate cartilage integrity, meniscal status, and ligamentous structures, ensuring the absence of significant arthritis in other compartments or patellofemoral joint.
  
  \item **Medication Review:** Patients are instructed to stop certain medications, such as blood thinners (e.g., aspirin, warfarin, novel oral anticoagulants) several days to a week prior to surgery, under the guidance of their prescribing physician. Other medications, like insulin or steroids, may require dosage adjustments.
  
  \item **NPO Status:** Patients must adhere to "nil per os" (NPO) guidelines, meaning no food or drink for a specified period (typically 6-8 hours) before surgery to prevent aspiration during anesthesia.
  
  \item **Preoperative Education:** Patients attend educational sessions to understand the procedure, potential risks, expected recovery, and rehabilitation protocol.
  
  \item **Informed Consent:** Detailed discussion with the surgeon regarding the procedure, alternatives, risks, and benefits, culminating in the signing of an informed consent form.
  
  \item **Antibiotics:** Prophylactic antibiotics are typically administered intravenously just before the incision to minimize the risk of infection.
\end{itemize}

This comprehensive preparation is vital for minimizing complications and ensuring a successful surgical outcome.

**Absolute Contraindications:**

\begin{itemize}
  \item Inflammatory arthritis (e.g., rheumatoid arthritis, psoriatic arthritis), as these are systemic diseases affecting the entire joint.
  \item Significant arthritis in more than one compartment of the knee (medial, lateral, or patellofemoral).
  \item Anterior or posterior cruciate ligament insufficiency, as these ligaments are critical for UKR stability.
  \item Fixed varus or valgus deformity exceeding 10-15 degrees, which cannot be corrected by soft tissue release.
  \item Active infection in the knee joint or elsewhere in the body.
  \item Neuropathic arthropathy (Charcot joint).
  \item Severe obesity (BMI > 35-40 kg/m\textasciicircum{}2), which increases stress on the implant and surgical risks.
  \item Insufficient bone stock or severe osteoporosis, compromising implant fixation.
  \item Previous knee infection or septic arthritis.
\end{itemize}

**Relative Contraindications \& Precautions:**

\begin{itemize}
  \item Younger age (under 55-60 years) is a relative concern due to potential for longer implant lifespan requirements, though active younger patients with isolated arthritis can be good candidates.
  \item High-demand activities or occupations that place excessive stress on the knee.
  \item Significant patellofemoral pain or early patellofemoral arthritis, which may progress.
  \item Moderate to severe osteopenia.
  \item Morbid obesity (BMI > 40 kg/m\textasciicircum{}2).
  \item Unrealistic patient expectations regarding outcomes or activity levels.
  \item Peripheral vascular disease or uncontrolled diabetes, which can impair wound healing.
\end{itemize}

Understanding these contraindications is essential for appropriate patient selection and optimizing surgical outcomes.

\section*{Surgical Technique \& Procedure}
The unicompartmental knee replacement procedure is typically performed under general anesthesia or regional anesthesia (spinal or epidural), often combined with a femoral nerve block for postoperative pain control. The patient is positioned supine on the operating table, and the affected leg is prepped and draped in a sterile fashion.

The surgical approach usually involves a smaller incision, typically 6-10 cm, on the anterior aspect of the knee, centered over the affected compartment. The quadriceps mechanism is carefully retracted, and the joint capsule is opened to expose the diseased compartment. Specialized instruments are used to precisely resect the damaged articular cartilage and a minimal amount of underlying bone from the femoral condyle and tibial plateau. Care is taken to preserve the anterior and posterior cruciate ligaments, as well as the menisci in the unaffected compartment. Trial components are then inserted to assess joint stability, alignment, and range of motion, ensuring proper soft tissue tension. Once optimal fit and function are confirmed, the definitive prosthetic components are implanted. These typically consist of a metal femoral component, a metal tibial baseplate, and a polyethylene bearing insert. The components are usually fixed to the bone using bone cement, though cementless options exist. After thorough irrigation, the joint capsule and soft tissues are meticulously repaired layer by layer, and the skin incision is closed with sutures or staples. A sterile dressing is applied.

Key operative steps include:

\begin{itemize}
  \item Anesthesia administration and patient positioning.
  \item Surgical incision and exposure of the affected knee compartment.
  \item Precise bone cuts on the femoral condyle and tibial plateau using specialized jigs.
  \item Insertion of trial components to assess fit, alignment, and ligamentous balance.
  \item Implantation of the definitive metal femoral component, metal tibial baseplate, and polyethylene bearing.
  \item Cementation of components (if using cemented implants).
  \item Thorough irrigation and debridement of the surgical site.
  \item Layered closure of the joint capsule, fascia, subcutaneous tissue, and skin.
\end{itemize}

This meticulous approach ensures optimal outcomes and minimizes complications.

\section*{Complications \& Risks}
As with any surgical procedure, unicompartmental knee replacement carries potential risks, which can be broadly categorized:

\begin{itemize}
  \item **General Surgical Risks:**
    \begin{itemize}
      \item **Anesthesia Complications:** Reactions to anesthesia, respiratory problems, cardiac events.
      \item **Bleeding:** Intraoperative or postoperative hemorrhage, potentially requiring blood transfusion.
      \item **Infection:** Superficial wound infection or deep joint infection, which can be devastating and may require further surgery, including implant removal.
      \item **Blood Clots:** Deep vein thrombosis (DVT) in the leg or pulmonary embolism (PE) in the lungs, potentially life-threatening.
    \end{itemize}
  
  \item **Procedure-Specific Risks:**
    \begin{itemize}
      \item **Implant Loosening or Failure:** The prosthetic components may loosen from the bone over time, requiring revision surgery.
      \item **Polyethylene Wear:** The plastic insert can wear out, necessitating replacement.
      \item **Progression of Arthritis:** Arthritis may develop or worsen in the previously healthy compartments, eventually requiring conversion to a total knee replacement.
      \item **Stiffness (Arthrofibrosis):** Limited range of motion due to scar tissue formation.
      \item **Nerve or Vessel Damage:** Injury to nerves (e.g., peroneal nerve) causing numbness or weakness, or damage to blood vessels.
      \item **Fracture:** Periprosthetic fracture during or after surgery.
      \item **Instability:** Ligamentous imbalance or improper implant positioning leading to a feeling of instability or dislocation.
      \item **Persistent Pain:** Despite successful surgery, some patients may experience ongoing pain.
      \item **Patellofemoral Pain:** New or worsening pain around the kneecap.
      \item **Component Malposition:** Incorrect alignment or sizing of the implants, affecting knee mechanics.
    \end{itemize}
\end{itemize}

Awareness of these complications is crucial for both the surgical team and the patient to ensure informed consent and appropriate postoperative management.

\section*{Postoperative Care \& Recovery}
Following unicompartmental knee replacement, the patient is closely monitored in the Post-Anesthesia Care Unit (PACU) for vital signs, pain levels, and surgical site integrity. Pain management is initiated, often involving a combination of oral medications, nerve blocks, and sometimes patient-controlled analgesia (PCA). Early mobilization is encouraged, typically within hours of surgery, with the assistance of physical therapists.

In the first 24-48 hours, the focus is on pain control, wound care, and initiating physical therapy. Patients are taught exercises to improve knee range of motion and strength. Most patients are discharged home within 1-3 days, depending on their recovery progress and pain control. For the first 1-2 weeks post-discharge, patients continue with prescribed pain medication, wound care, and a structured home exercise program or outpatient physical therapy. Swelling and bruising are common and can be managed with elevation and ice. The goal is to regain full extension and achieve at least 90 degrees of flexion. Long-term recovery involves continued physical therapy for several weeks to months, gradually increasing activity levels. Most patients can return to light activities and driving within 3-6 weeks, with full recovery and return to more demanding activities potentially taking 3-6 months.

During the unicompartmental knee replacement procedure, the surgeon documents the condition of the knee joint, noting the extent of cartilage degeneration, bone quality, and any associated pathologies such as meniscal tears or ligamentous injuries. The pathology report on any resected tissues typically confirms the diagnosis of osteoarthritis, detailing the degree of cartilage loss and any other relevant findings. This documentation is crucial for understanding the surgical outcomes and guiding future management.

A normal, successful postoperative course following unicompartmental knee replacement is characterized by significant pain resolution, improved knee function, and a gradual return to preoperative activities. Patients typically experience a reduction in pain levels within the first few weeks, with most achieving substantial pain relief by 6-12 weeks post-surgery. The timeline for healing varies, but many patients can expect to regain full range of motion and strength within 3-6 months, allowing them to return to light sports and recreational activities. Regular follow-up appointments are essential to monitor recovery and address any concerns that may arise.

Postoperative care is critical for ensuring optimal recovery following unicompartmental knee replacement. Patients are typically scheduled for follow-up visits at 2 weeks for wound check and suture/staple removal, and then at 6 weeks, 3 months, 6 months, and 1 year post-surgery. 

Key components of postoperative care include:

\begin{itemize}
  \item **Physical Therapy:** A structured rehabilitation program is essential, starting immediately after surgery and continuing for several weeks to months, focusing on range of motion, strength, and gait training.
  
  \item **Imaging:** Follow-up X-rays are usually taken at 6 weeks, 3 months, 1 year, and then periodically (e.g., every 2-5 years) to monitor implant position, alignment, and detect any signs of loosening or wear.
  
  \item **Activity Resumption:** Gradual return to activities is guided by the surgeon and physical therapist. Light activities can often resume within weeks, with more strenuous activities taking several months.
  
  \item **Warning Signs:** Patients are educated on warning signs to report to their doctor, such as increasing pain, redness, swelling, fever, drainage from the wound, or sudden instability.
\end{itemize}

Long-term annual clinical review is often recommended to monitor the overall health of the knee and the implant.

\section*{Outcomes \& Prognosis}
Unicompartmental knee replacement accounts for a smaller but growing proportion of all knee arthroplasties performed worldwide, typically ranging from 5\% to 15\% of total knee replacements. The incidence of UKR has been increasing due to improved patient selection, surgical techniques, and implant designs, leading to better long-term outcomes. It is most commonly performed in patients over the age of 55, though younger, active individuals with isolated unicompartmental osteoarthritis are increasingly considered candidates. There is no significant sex predilection, with incidence rates being relatively similar between men and women. The most frequent indication is medial compartment osteoarthritis, often associated with a varus deformity. While the overall mortality rate associated with UKR is very low, similar to TKR, the morbidity profile is generally lower, with fewer complications such as blood loss, transfusion rates, and length of hospital stay compared to total knee replacement.

The outcomes and prognosis for unicompartmental knee replacement are generally excellent for carefully selected patients, with reported implant survival rates comparable to total knee replacement in the long term. Success rates, defined by pain relief and functional improvement, are often in the range of 85-95\% at 10 years and 70-80\% at 15-20 years. Patients typically experience significant pain reduction, improved range of motion, and a more natural-feeling knee compared to TKR, allowing for a higher level of activity. Functional outcomes, as measured by patient-reported outcome measures (PROMs) like the Oxford Knee Score or WOMAC, consistently show substantial improvement. The main reasons for revision include progression of arthritis in other compartments, aseptic loosening of components, and polyethylene wear. Prognostic factors for good outcomes include appropriate patient selection (isolated arthritis, intact ligaments), experienced surgical teams, and adherence to postoperative rehabilitation. While UKR has a higher revision rate than TKR, the revision surgery to a total knee replacement is generally straightforward and yields good results.

\section*{References}
\begin{enumerate}
  \item MedlinePlus. Partial knee replacement. U.S. National Library of Medicine; 2025. Available from: \url{https://medlineplus.gov/ency/article/007238.htm}
  
  \item Scott WN, Insall JN, Windsor RE, et al. Insall \& Scott Surgery of the Knee. 6th ed. Elsevier; 2025.
  
  \item American Academy of Orthopaedic Surgeons (AAOS). Unicompartmental Knee Arthroplasty. AAOS Clinical Practice Guideline; 2024.
  
  \item Ranawat AS, et al. Unicompartmental Knee Arthroplasty: Indications, Surgical Technique, and Outcomes. J Am Acad Orthop Surg. 2023;31(1):e1-e12.
\end{enumerate}

\clearpage
